Fix \((\mu,\mathcal M)\in\mathfrak M_{m,\epsilon}\).  All conditional equalities below hold for \(\mu\)-almost every \(q\in K_{\epsilon}\) at which the stipulated regular conditional laws are defined.  Put \(r=p(q)=q_{10}+q_{11}\).  By observed-cell calibration,
\[
 \Pr(A_{ij}=1\mid Q_i=q)=r,
 \qquad
 \Pr(A_{ij}=0\mid Q_i=q)=1-r.
\]
Since \(q\in K_{\epsilon}\), \(r\geq\epsilon>0\) and \(1-r\geq\epsilon>0\), so every division below is valid.  Latent ignorability gives, for \(a\in\{0,1\}\),
\[
 \mathbb E\{Y_{ij}(a)\mid Q_i=q\}
 =\mathbb E\{Y_{ij}(a)\mid A_{ij}=a,Q_i=q\}.
\]
On \(A_{ij}=a\), consistency gives \(Y_{ij}=Y_{ij}(a)\).  Calibration and division by the positive treatment probabilities therefore yield
\[
 \begin{aligned}
 \mathbb E\{Y_{ij}(1)\mid Q_i=q\}
 &=\Pr(Y_{ij}=1\mid A_{ij}=1,Q_i=q)=\frac{q_{11}}r,\\
 \mathbb E\{Y_{ij}(0)\mid Q_i=q\}
 &=\Pr(Y_{ij}=1\mid A_{ij}=0,Q_i=q)=\frac{q_{01}}{1-r}.
 \end{aligned}
\]
Subtracting and using the definition of \(f\) proves
\[
 \mathbb E\{Y_{ij}(1)-Y_{ij}(0)\mid Q_i=q\}=f(q).
\]
The environment-law assumption gives \(Q_i\sim\mu\).  Binary potential outcomes are integrable, so iterated expectation gives
\[
 \mathbb E\{Y_{ij}(1)-Y_{ij}(0)\}
 =\int_{K_{\epsilon}}\left\{\frac{q_{11}}{p(q)}-\frac{q_{01}}{1-p(q)}\right\}\,d\mu(q)
 =\int f\,d\mu=\tau(\mu).
\]
This is the cluster-equally-weighted ATE because all clusters have the common fixed plate size \(m\), clusters are identically distributed by independent cluster sampling, and individuals are conditionally identically distributed by the conditional iid plate assumption.

Conditional on \(Q_i=q\), the causal vectors are mutually independent over \(j\) by the conditional iid plate assumption.  Each \(O_{ij}=(A_{ij},Y_{ij})\) is a measurable function of its causal vector by consistency, so the observed pairs are conditionally iid.  Calibration gives their cell probabilities \(q\).  Thus, for every \(c\in\mathcal I_m\),
\[
 \Pr(C_i=c\mid Q_i=q)
 =\frac{m!}{\prod_{a,y}c_{ay}!}\prod_{a,y}q_{ay}^{c_{ay}}
 =B_c^{(m)}(q).
\]
Integrating with respect to \(Q_i\sim\mu\) gives
\[
 \Pr(C_i=c)=\int_{K_{\epsilon}}B_c^{(m)}(q)\,d\mu(q).
\]

Conversely, fix \(\mu\in\mathcal P(K_{\epsilon})\).  Independently over clusters draw \(Q_i\sim\mu\).  Given \(Q_i=q\), put \(r=p(q)\) and, independently over \(j\), draw mutually independent variables
\[
 A_{ij}\sim\operatorname{Bernoulli}(r),\qquad
 Y_{ij}(0)\sim\operatorname{Bernoulli}\!\left(\frac{q_{01}}{1-r}\right),\qquad
 Y_{ij}(1)\sim\operatorname{Bernoulli}\!\left(\frac{q_{11}}r\right).
\]
Set \(Y_{ij}=A_{ij}Y_{ij}(1)+(1-A_{ij})Y_{ij}(0)\).  Overlap makes both denominators positive.  Since \(q\in\Delta_3\),
\[
 0\leq q_{11}\leq r,
 \qquad
 0\leq q_{01}\leq1-r,
\]
so both Bernoulli means lie in \([0,1]\).  The construction is well defined.  Its conditional mutual independence proves latent ignorability; independent repetitions prove the conditional iid plate condition; the definition of \(Y_{ij}\) proves consistency; and independent cluster draws prove independent identically distributed cluster sampling and the environment-law condition.  Finally,
\[
 \begin{array}{ll}
 \Pr(A_{ij}=1,Y_{ij}=1\mid q)=r(q_{11}/r)=q_{11},
 &\Pr(A_{ij}=1,Y_{ij}=0\mid q)=r-q_{11}=q_{10},\\[4pt]
 \Pr(A_{ij}=0,Y_{ij}=1\mid q)=(1-r)(q_{01}/(1-r))=q_{01},
 &\Pr(A_{ij}=0,Y_{ij}=0\mid q)=1-r-q_{01}=q_{00}.
 \end{array}
\]
Thus calibration holds and every defining property of \(\mathfrak M_{m,\epsilon}\) is satisfied.  The forward calculation applied to this construction gives its multinomial count law and causal target.